\documentclass[10pt,a4paper]{amsart}
\usepackage[margin=19mm]{geometry}
\usepackage{amsmath,amssymb,mathtools}
\usepackage[scr=rsfs,cal=euler]{mathalfa}
\usepackage{xfrac}
\usepackage[unicode,hidelinks,bookmarks=false]{hyperref}
\newcommand{\ie}{i.e.\ }
\newcommand{\cf}{cf.\ }
\newcommand{\resp}{respectively\ }
\newcommand{\Subsection}{\subsection}
\providecommand{\nobreakdash}{\nobreak\hbox{-}}
\setlength{\emergencystretch}{2em}
\multlinegap0pt
\usepackage{bm}
\usepackage{bbm}
\usepackage{dsfont}
\usepackage{xspace}
\usepackage{cancel}
\usepackage{mathtools}
\usepackage{amsthm}
\makeatletter
\@ifundefined{Thm}{\newtheorem{Thm}{Thm}}{}
\@ifundefined{Cor}{\newtheorem{Cor}[Thm]{Corollary}}{}
\@ifundefined{Lem}{\newtheorem{Lem}[Thm]{Lemma}}{}
\@ifundefined{Prop}{\newtheorem{Prop}[Thm]{Proposition}}{}
\makeatother
\newcommand{\connect}{\xleftrightarrow}
\title[An alternative approach for the mean-field behaviour of spre]{An alternative approach for the mean-field behaviour of spread-out Bernoulli percolation in dimensions $d>6$}
\author{Hugo Duminil-Copin, Romain Panis}
\date{Source: arXiv:2410.03647v2 (CC BY 4.0)}
\begin{document}
\maketitle

\noindent\emph{Source and licence.} An alternative approach for the mean-field behaviour of spread-out Bernoulli percolation in dimensions $d>6$. Version arXiv:2410.03647v2 (CC BY 4.0), distributed under the Creative Commons Attribution 4.0 International licence, as recorded in the arXiv licence metadata for this exact version and shown on the abstract page https://arxiv.org/abs/2410.03647v2. Full rights notice: licenses/arxiv-2410.03647-CC-BY-4.0.txt.\par\smallskip

\noindent\textit{Editorial scope.} Counted unit: the Prop whose source label is \texttt{prop:regularity perco 2} in the article (source0.tex lines 1413--1418, source label \texttt{prop:regularity perco 2}), delivered together with the article's own complete original proof of it (source0.tex lines 1442--1503, one \texttt{proof} environment of 62 lines). No mathematics is written, repaired or re-derived: statement and proof are the article's own text line for line. The proof is detached from the statement in the source: the article states the result at lines 1413--1418 and prints its proof at lines 1442--1503, and the proof environment's own header names the statement, so the association is the article's own. Required context retained uncounted: Lem: SL upper bound (lines 235--239); Lem: SL lower bound (lines 363--377); prop: final prop proof mainperco thm (lines 1195--1213); prop:regularitygeneralperco (lines 1314--1324); prop: exponential moment exit time (lines 2002--2007); cor: estimate srw (lines 2034--2038); prop:uniform Harnack (lines 2064--2069). Every definition, display and label of those lines is retained unchanged. Omitted from this package and not redistributed: 1--234, 240--362, 378--1194, 1214--1313, 1325--1412, 1419--1441, 1504--2001, 2008--2033, 2039--2063, 2070--2343. Nothing inside a retained excerpt is omitted or altered: every retained body is delivered exactly as the article's lines are written, so no omission receipt of any kind accompanies this package. Citations: the retained excerpts carry no citation command at all, so no bibliography entry of the article is transcribed and no reference is redistributed. Reference adaptation: none was needed and none was made; every reference inside the delivered unit resolves inside it, so no unresolved reference or citation is produced. Printed slips kept literal: no printed word, symbol, hypothesis or number of the counted statement or of its proof was altered. Layout and class: the article's own class is \texttt{article} with packages including amsfonts, amsthm, amsmath, amssymb, amscd, mathrsfs, eucal, graphicx, epstopdf, pict2e, epic, xcolor, mathtools, float; the wrapper is the lane's pinned \texttt{amsart} editorial wrapper at 10pt on A4, which restates the theorem-like environments the retained text needs and adds the article's own macro definitions that the retained text needs (\texttt{\textbackslash connect}), with extra wrapper packages (bm, bbm, dsfont, xspace, cancel, mathtools). The delivered numbering is the unit's own. Assets: no retained span contains a figure environment, a graphics-inclusion command, a PDF-page inclusion, a table or a TikZ picture; no image file is copied into this package, the package declares no asset dependency and no image file is redistributed with it. Licence: version arXiv:2410.03647v2 of this article is distributed under the Creative Commons Attribution 4.0 International licence, as recorded in the arXiv licence metadata for this exact version and shown on the abstract page https://arxiv.org/abs/2410.03647v2. A full rights notice is delivered at licenses/arxiv-2410.03647-CC-BY-4.0.txt. Characters: every retained excerpt is pure ASCII, so the delivered engine, XeLaTeX with its default Latin Modern text font, reports no missing character.
\begin{Lem}\label{Lem: SL upper bound} Let $d\geq 2$, $\beta>0$, $o\in S\subset \Lambda$, and $x\in \Lambda$. Then,
\begin{equation}\label{eq:SLperco}
	\mathbb P_{\beta}[o\connect{\Lambda\:}x]\le \mathbb P_{\beta}[o\connect{S\:}x]+ \sum_{\substack{y\in S\\ z\in \Lambda\setminus S}}\mathbb P_{\beta}[o\connect{S\:}y]p_{yz,\beta}\mathbb P_{\beta}[z\connect{\Lambda\:}x].
\end{equation}
\end{Lem}

\begin{Lem}\label{Lem: SL lower bound} Let $d\geq 2$, $0<\beta<\beta_c$, $o\in S\subset \Lambda$, and $x\in \Lambda$. Then,
\begin{equation}\label{eq:reversed SL perco}
\mathbb P_{\beta}[o\connect{\Lambda\:}x]\ge \mathbb P_{\beta}[o\connect{S\:}x]+\sum_{\substack{y\in S\\ z\in \Lambda\setminus S}}\mathbb P_{\beta}[o\connect{S\:}y]p_{yz,\beta} \mathbb P_{\beta}[z\connect{\Lambda\:}x]-E_\beta(S,\Lambda,o,x),
\end{equation}
where
\begin{align}\label{eq:error first term}
&E_\beta (S,\Lambda,o,x)
:=\sum_{u,v\in S}\sum_{\substack{y \in S\\ z\in\Lambda\setminus S}}\mathbb P_{\beta}[o\connect{S\:}u]\mathbb P_{\beta}[u\connect{S\:}y]p_{yz,\beta}\mathbb P_{\beta}[z\connect{\Lambda\:}v]\mathbb P_{\beta}[u\connect{S\:}v]\mathbb P_\beta[v\connect{\Lambda\:}x]
\\\label{eq:error second term}
&+\sum_{\substack{u\in S\\ v\in \Lambda}}\sum_{\substack{y,s\in S\\ z,t\in \Lambda\setminus S\\ yz\neq st}}\mathbb P_\beta[o\connect{S\:}u]\mathbb P_\beta[u\connect{S\:}y]\mathbb P_\beta[u\connect{S\:}s]p_{yz,\beta}p_{st,\beta}\mathbb P_\beta[z\connect{\Lambda\:}v]\mathbb P_\beta[t\connect{\Lambda\:}v]\mathbb P_\beta[v\connect{\Lambda\:}x].
%\\\label{eq:error third term}
%&+\sum_{\substack{y, s\in S\\y\neq s\\ z,t\in \Lambda \setminus S\\ z\neq t}}\mathbb P_\beta[o\connect{S\:}u]\mathbb P_\beta[u \connect{S\:}y]\mathbb P_\beta[u\connect{S\:}s]p_{yz,\beta}p_{st,\beta}\mathbb P_\beta[z\connect{\Lambda\:}v]\mathbb P_\beta[t\connect{\Lambda\:}v]
%	\\\label{eq:error fourth term}	&+ \delta_o(u)\sum_{\substack{y\in S\\ z,t\in \Lambda\setminus S\\ z\neq t}}\mathbb P_\beta[o\connect{S\:}y]p_{yz,\beta}p_{yt}(\beta)\mathbb P_\beta[z\connect{\Lambda\:}v]\mathbb P_\beta[t\connect{\Lambda\:}v].
\end{align}
\end{Lem}

\begin{Prop}\label{prop: final prop proof mainperco thm}
Fix $d>6$. There exist $C,L_0\geq 1$ such that for every $L\geq L_0$,
\begin{align}
\beta_c&\leq 1+\frac{C}{L^d}, \label{eq:bound final beta_c}
\\
\varphi_{\beta_c}(B)&\le 1+\frac{C}{L^d} &\forall B\in\mathcal B,\label{eq:bound final phi block}
%\\
%\varphi_{\beta_c}(\Lambda_n)&\le 1+\frac{K}{L^d} &\forall n\ge0,\label{eq:bound final phi(Lambda_n)}
\\
E_{\beta_c}(B)&\leq \frac{C}{L^d} &\forall B\in \mathcal B,\label{eq:error amplitude general blocks}
\\
\psi_{\beta_c}(\mathbb H_n)&\le \frac{C}{L} &\forall n\ge 0, \label{eq:bound final psi}
\\
\mathbb P_{\beta_c}[0\connect{}x] &\le \frac{C}{L^d}\left(\frac{L}{L\vee |x|}\right)^{d-2}&\forall x\in (\mathbb Z^d)^*,\label{eq:upper bound below L perco}
\\
\mathbb P_{\beta_c}[0\connect{\mathbb H\:}x]&\le\frac{C}{L^d}\left(\frac{L}{L\vee |x_1|}\right)^{d-1}&\forall x\in \mathbb H^*\label{eq:upper half-spacebound below L perco}.
%E^{\mathbb H_n,\mathbb Z^d}_\beta&\leq \frac{K}{L^d} &\forall n\geq 0,\label{eq:error amplitude half-space}
\end{align} 
\end{Prop}

\begin{Prop}[Regularity estimate at mesoscopic scales]\label{prop:regularitygeneralperco} For every $\eta>0$, there exist $\delta=\delta(\eta,d)\in(0,1/2)$, and $L_5=L_5(\eta,d)\geq L_0$ such that for every $L\geq L_5$, every $\beta\leq \beta_c$, every $n\geq \delta^{-1}L$, every $\Lambda\supset \Lambda_{3n}$, every $X\subset \Lambda\setminus\Lambda_{3n}$, and every $u,v\in \Lambda_{\lfloor \delta n\rfloor}$,
\begin{multline}\label{eq:reg mesoscopic boosted}
\left|\sum_{x\in X}\Big(\mathbb P_\beta[u\connect{\Lambda\:}x]-\mathbb P_\beta[v\connect{\Lambda\:}x]\Big)\right|\le \eta\max_{w\in \Lambda_{3n}}
\sum_{x\in X}\mathbb P_\beta[w\connect{\Lambda\:}x]\\+\delta^{-1}\max_{w\in \Lambda_{3n}}\max_{\substack{S\in \mathcal B\\S+w\subset \Lambda_{3n}}}\sum_{x\in X}E_\beta(S+w,\Lambda,w,x).
\end{multline}
Moreover, if $u,v \in \Lambda_n$ are such that $|u-v|\leq 2L$, one has
\begin{equation}\label{eq:uv close in the statement}
	\left|\sum_{x\in X}\Big(\mathbb P_\beta[u\connect{\Lambda\:}x]-\mathbb P_\beta[v\connect{\Lambda\:}x]\Big)\right|\le \eta\max_{w\in \Lambda_{3n}}
\sum_{x\in X}\mathbb P_\beta[w\connect{\Lambda\:}x]\\+\delta^{-1}\max_{w\in \Lambda_{3n}}\sum_{x\in X}E_\beta(\{w\},\Lambda,w,x).
\end{equation}
\end{Prop}

\begin{Prop}\label{prop: exponential moment exit time} Let $d\geq 1$. For every $n,m\geq 1$ with $n\geq m$, every $\mu \in \mathcal P_m$, and every $u\in \Lambda_n$,
\begin{equation}
	\mathbb E_{\mu,u}[\tau_n]\leq 9d\left(\frac{n}{m}\right)^2.
\end{equation}

\end{Prop}

\begin{Cor}\label{cor: estimate srw} Let $d\geq 1$ and $A,\eta>0$. Let $n,m\geq 1$ with $(n/m)\leq A$. There exist $T,\varepsilon>0$ which depend on $(A, \eta,d)$ such that the following holds. For every $\mu \in \mathcal P_m$, every $\varphi\in [1-\varepsilon,1+\varepsilon]$, every $f:\mathbb Z^d\rightarrow \mathbb R^+$, and every $u \in \Lambda_n$,
\begin{equation}
	\Big|\mathbb E_{\mu,u}[\varphi^{\tau_n\wedge T}f(X_{\tau_n\wedge T})]-\mathbb E_{\mu,u}[f(X_{\tau_n})]\Big|\leq \eta \max\{f(w):w\in \Lambda_{n+2m}\}.
\end{equation}
\end{Cor}

\begin{Prop}[Uniform Harnack-type estimate]\label{prop:uniform Harnack} Let $d\geq 3$, $\alpha>0$, and $\eta>0$. There exists $C_{\rm RW}=C_{\rm RW}(\alpha,d)>0$ and $N_1=N_1(\eta,\alpha,d)>0$ such that the following holds. For every $n,m\geq 1$ satisfying $\tfrac{n}{m}\geq N_1$, every $\mu\in \mathcal P_m$, every $f:\mathbb Z^d\rightarrow \mathbb R^+$, and every $u,v\in \Lambda_n$,
\begin{multline}
	\mathbb E_{\mu,u}[f(X_\tau)]\leq C_{\rm RW}\mathbb E_{\mu,v}[f(X_\tau)] +\eta\max\{f(w):w\in \Lambda_{(1+\alpha)n+2m}\}\\+2C_{\rm RW}\max\Big\{|f(w)-f(w')|: w,w'\in \Lambda_{3m(n/m)^{1/10}}(z), \: z\in \partial \Lambda_{(1+\alpha)n}\Big\},
\end{multline}
where $\tau:=\inf \{k\geq 0: X_k\notin \Lambda_{(1+\alpha)n}\}$.
\end{Prop}

\begin{Prop}[Regularity estimate at macroscopic scales]\label{prop:regularity perco 2} Let $\alpha>0$. There exists $C_{\rm RW}=C_{\rm RW}(\alpha)>0$, and for every $\eta>0$, there exist $A=A(\alpha,\eta)>0$, $L_6\geq L_0$, and $\varepsilon_0=\varepsilon_0(\eta,\alpha)>0$ such that the following holds. For every $L\geq L_6$, every $\varepsilon<\varepsilon_0$, every $n\le 6L_\beta(\varepsilon)$  with $n\geq AL$, every $\beta\le \beta_c$, every $\Lambda\supset\Lambda_{(1+\alpha)n}$, and every $x\in \Lambda\setminus \Lambda_{(1+\alpha)n}$,
\begin{align}
\max_{w\in \Lambda_n}\mathbb P_\beta[w\connect{\Lambda\:}x]&\le C_{\rm RW}\min_{w\in \Lambda_n}\mathbb P_\beta[w\connect{\Lambda\:}x]+\eta \max_{w\in \Lambda_{(1+\alpha)n}}\mathbb P_\beta[w\connect{\Lambda\:}x]\notag
\\&\quad+A\max_{w\in \Lambda_{(1+\alpha)n}}\max_{\substack{S\in \mathcal B\\S+w\subset \Lambda_{(1+\alpha)n}}}{E}_\beta(S+w, \Lambda,w,x).
\end{align}
\end{Prop}

\begin{proof}[Proof of Proposition \textup{\ref{prop:regularity perco 2}}] Let $\alpha,\eta>0$. Set $M:=\lfloor \tfrac{\alpha}{6} n\rfloor$, and let $C_{\rm RW}$ and $N_1$ be given by Proposition \ref{prop:uniform Harnack} applied to $\tfrac{\alpha}{6}$ and $\tfrac{\eta}{4}$.
Let $\delta$ and $L_5$ be given by Proposition \ref{prop:regularitygeneralperco} with $\frac{\eta}{4C_{\rm RW}}$ and let $L\geq L_5$. To the cost of diminishing $\delta$, we additionally assume that
\begin{equation}\label{eq:take delta smaller}
	\delta\leq (N_1\alpha)^{-1}.
\end{equation}
Recall the random walk distribution $\mathbb P^{{\rm RW},m}_u$ introduced above. Let $\tau$ be the hitting time of $\mathbb Z^d\setminus\Lambda_{n+M}$. If $m=\lfloor (\delta\alpha/36)^{10/9}n\rfloor$, then $n/m\geq N_1$ by \eqref{eq:take delta smaller}. We further assume that $n\geq (\delta\alpha/36)^{-10/9}(L+1)$ so that $m\geq L$ and Proposition \ref{prop:uniform Harnack} applies. Using this latter result, we get, for every $u,v\in \Lambda_n$,
\begin{multline}\label{eq:proof propreg2 1}
	\mathbb E^{{\rm RW},m}_u\Big[\mathbb P_\beta[X_\tau \connect{\Lambda\:}x]\Big]\leq C_{\rm RW}\mathbb E^{{\rm RW},m}_v\Big[\mathbb P_\beta[X_\tau \connect{\Lambda\:}x]\Big]+\frac{\eta}{4}\max\{\mathbb P_\beta[w \connect{\Lambda\:}x]:w\in \Lambda_{(1+\alpha)n}\}\\+2C_{\rm RW}\max \Big\{|\mathbb P_\beta[w \connect{\Lambda\:}x]-\mathbb P_\beta[w'\connect{\Lambda\:}x]|: w,w'\in \Lambda_{\lfloor \delta M\rfloor}(z), \: z\in \partial \Lambda_{n+M}\Big\} 
\end{multline}
where we used that, for this choice of $M,m$ and for $n$ large enough, one has: 
\begin{align}
3m(n/m)^{1/10}&\leq \frac{\delta \alpha}{12}n\leq \frac{\delta\alpha}{6}n-\delta-1\leq \lfloor\delta M\rfloor,\\
n+M+2m&\leq (1+\alpha)n.\end{align} 
Combining \eqref{eq:proof propreg2 1} with Corollary \ref{cor: estimate srw} (applied to $\frac{\eta}{8C_{\rm RW}}$ and $A\approx (\alpha\delta)^{-10/9}$) provides $T=T(\eta/(8C_{\rm RW}),\alpha,d)>0$ and $\varepsilon_0=\varepsilon_0(T)$ such that $(1+\varepsilon_0)^T\leq 2$, and for every $\varphi\in [1-\varepsilon_0,1+\varepsilon_0]$,
\begin{multline}\label{eq:proof propreg2 1.5}
	\mathbb E_u^{{\rm RW},m}\Big[\varphi^{\tau\wedge T}\mathbb P_\beta[X_{\tau\wedge T}\connect{\Lambda\:}x]\Big]\leq C_{\rm RW}\mathbb E_v^{{\rm RW},m}\Big[\varphi^{\tau\wedge T}\mathbb P_\beta[X_{\tau\wedge T}\connect{\Lambda\:}x]\Big]\\+\frac{\eta}{2}\max\{\mathbb P_\beta[w\connect{\Lambda\:}x]:w\in \Lambda_{(1+\alpha)n}\}\\+2C_{\rm RW}\max \Big\{|\mathbb P_\beta[w\connect{\Lambda\:}x]-\mathbb P_\beta[w'\connect{\Lambda\:}x]|: w,w'\in \Lambda_{\lfloor \delta M\rfloor}(z), \: z\in \partial \Lambda_{n+M}\Big\}.
\end{multline}
Thanks to Proposition \ref{prop: exponential moment exit time}, for every $u\in \Lambda_n$,
\begin{equation}\label{eq:proof propreg2 2}
	\mathbb E_u^{{\rm RW},m}[\tau]\leq C_1(1+\alpha)^2\delta^{-20/9},
\end{equation}
for some $C_1=C_1(d)>0$.
%
% We will use two a priori estimates on the random walk and the stopping time $\tau$, that can be easily obtained from classical random walk analysis. Indeed, applying Proposition \ref{prop: estimates non srw} to $\mu$, $m$, and $\gamma:=8\delta^{-1}(1+\tfrac{\alpha}{6})$ (i.e. $n+M=\gamma m$), there exists $C_0$ which only depend on $d$ such that 
% \begin{align}\label{eq:hb}
%\mathbb E_u'[\tau]&\le C_0 \gamma^2 &\forall u\in\Lambda_n,\\
%\mathbb P_{u}'[X_\tau'\in B_i]&\leq C_0\mathbb P_v'[X_\tau'\in B_i] &\forall u,v\in \Lambda_n,\:\forall 1\le i\le q.\label{eq:hc}
%\end{align}
%We fix $T=T(\eta,\alpha,d)>0$ be large enough so that, by Markov's inequality and \eqref{eq:proof propreg2 2}, for every $u \in \Lambda_n$,
%\begin{equation}
%	\mathbb P_u'[\tau >T]\leq \frac{\mathbb E_u'[\tau]}{T}\leq \frac{\eta}{16}.
%\end{equation}
We are now in a position to prove the desired result.

Let $L\geq L_5$ and $\varepsilon<\varepsilon_0$. By choosing $L$ large enough we may additionally assume that $C/L^d\leq \varepsilon_0$ where we recall that $C$ is given by Proposition \ref{prop: final prop proof mainperco thm}. Fix $u,v\in \Lambda_n$. By definition of $m$ and \eqref{eq:take delta smaller}, we observe that $6m<n$. Since we also have $n\le 6L_\beta(\varepsilon)$, Proposition~\ref{prop: final prop proof mainperco thm} gives
\begin{equation}
1-\varepsilon\le \varphi_\beta(\Lambda_{m-1})\le 1+\frac{C}{L^d}\le 1+\varepsilon_0.
\end{equation}
We now set $\varphi:=\varphi_\beta(\Lambda_{m-1})$. Iterating the two bounds of Lemmata \ref{Lem: SL upper bound} and \ref{Lem: SL lower bound} with $S$ being translates of $\Lambda_{m-1}$, and until time $\tau\wedge T$ gives \begin{align}
\mathbb P_\beta[u\connect{\Lambda\:}x]&\le  \mathbb E_{u}^{{\rm RW},m}\Big[\varphi^{\tau\wedge T} \mathbb P_\beta[Y_{\tau\wedge T}\connect{\Lambda\:}x]\Big],\label{eq:hea perco}\\
\mathbb P_\beta[v\connect{\Lambda\:}x]&\ge \mathbb E_{v}^{{\rm RW},m}\Big[\varphi^{\tau\wedge T} \mathbb P_\beta[Y_{\tau\wedge T}\connect{\Lambda\:}x]\Big]-\mathbb E_{v}^{{\rm RW},m}\left[\sum_{t=0}^{\tau\wedge T-1}\varphi^t{E}_\beta(\Lambda_{m-1}(Y_t),\Lambda,Y_t,x)\right].\label{eq:hda perco}
\end{align}
Using that $(1+\varepsilon_0)^T\leq 2$ and \eqref{eq:proof propreg2 2}, we get
\begin{multline}\label{eq:proof propref2 2.5}
	\mathbb E_{v}^{{\rm RW},m}\left[\sum_{t=0}^{\tau\wedge T-1}\varphi^t{E}_\beta(\Lambda_{m-1}(Y_t),\Lambda,Y_t,x)\right]\\\leq 2C_1(1+\alpha)^2\delta^{-20/9}\max_{w\in \Lambda_{(1+\alpha)n}}\max_{\substack{S\in \mathcal B\\S+w\subset \Lambda_{(1+\alpha)n}}}{E}_\beta(S+w, \Lambda,w,x).
\end{multline}
Combining \eqref{eq:hea perco} and \eqref{eq:proof propreg2 1.5} gives
\begin{multline}\label{eq:proof propreg2 3}
	\mathbb P_\beta[u\connect{\Lambda\:}x]\leq C_{\rm RW}\mathbb E_v^{{\rm RW},m}\Big[\varphi^{\tau\wedge T}\mathbb P_\beta[Y_{\tau \wedge T}\connect{\Lambda\:}x]\Big]+\frac{\eta}{2}\max\{\mathbb P_\beta[w\connect{\Lambda\:}x]:w\in \Lambda_{(1+\alpha)n}\}\\+2C_{\rm RW}\max \Big\{|\mathbb P_\beta[w\connect{\Lambda\:}x]-\mathbb P_\beta[w'\connect{\Lambda\:}x]|: w,w'\in \Lambda_{\lfloor \delta M\rfloor}(z), \: z\in \partial \Lambda_{n+M}\Big\}. 
\end{multline}
Assume that $M\geq \delta^{-1}L$, which occurs as soon as $n\geq 6\alpha^{-1}\delta^{-1}(L+1)$. By Proposition \ref{prop:regularitygeneralperco} (recall that it is applied to $\tfrac{\eta}{4C_{\rm RW}}$),
\begin{multline}\label{eq:proof propreg2 4}
	2C_{\rm RW}\max \Big\{|\mathbb P_\beta[w\connect{\Lambda\:}x]-\mathbb P_\beta[w'\connect{\Lambda\:}x]|: w,w'\in \Lambda_{\lfloor \delta M\rfloor}(z), \: z\in \partial \Lambda_{n+M}\Big\}\\\leq \frac{\eta}{2}\max\{\mathbb P_\beta[w\connect{\Lambda\:}x]: w \in \Lambda_{n+3M}\}+2C_{\rm RW}\delta^{-1}\max_{w\in \Lambda_{(1+\alpha)n}}\max_{\substack{S\in \mathcal B\\S+w\subset \Lambda_{(1+\alpha)n}}}{E}_\beta(S+w, \Lambda,w,x).
\end{multline}
Set $B:=2C_{\rm RW}(C_1(1+\alpha)^2\delta^{-20/9}+\delta^{-1})$. 
Plugging \eqref{eq:hda perco}, \eqref{eq:proof propref2 2.5}, and \eqref{eq:proof propreg2 4} in \eqref{eq:proof propreg2 3} gives that, for every $n\geq \Big((\delta\alpha/36)^{-10/9}\vee 6\alpha^{-1}\delta^{-1}\Big)(L+1)$,
\begin{multline}
	\mathbb P_\beta[u\connect{\Lambda\:}x]\leq C_{\rm RW}\mathbb P_\beta[v\connect{\Lambda\:}x]+\eta\max\{\mathbb P_\beta[w\connect{\Lambda\:}x]:w\in \Lambda_{(1+\alpha)n}\}
	\\+B\max_{w\in \Lambda_{(1+\alpha)n}}\max_{\substack{S\in \mathcal B\\S+w\subset \Lambda_{(1+\alpha)n}}}{E}_\beta(S+w, \Lambda,w,x),
\end{multline}
where we used that $n+3M\leq (1+\alpha)n$. Since $u$ and $v$ are arbitrary in $\Lambda_n$, the proof follows readily from setting $A:=B\vee (\delta\alpha/36)^{-10/9}\vee 6\alpha^{-1}\delta^{-1}$.
 \end{proof}

\end{document}
